% =============================================================
%  Section: Conclusions
% =============================================================

This work reproduces the YOLOv1 detector of \citet{redmon2016you} on the COCO
2017 dataset, using a ResNet18 backbone~\citep{he2016deep} pretrained on
ImageNet, a convolutional adapter and the paper's two-layer dense detection
head. The pipeline---dataset preparation, model architecture, five-term
loss, distributed training loop with mixed precision, and mAP evaluation
with per-class non-maximum suppression---is implemented in PyTorch and is
end-to-end reproducible from the seeded dataset split.

The headline number is $\mathrm{mAP}@0.5 = 0.0415$ on the held-out COCO test
split (V6 checkpoint, prob\_threshold~$0.1$), with forward-pass latency of
$\sim\!5.3$~ms per image at $448\!\times\!448$ on a single Ada-generation
laptop GPU, well inside the real-time regime the original paper sought.
Both the absolute mAP and the gap to the paper's $63.4$ on VOC are
attributable to deliberate choices made within the budget of the project
rather than to implementation defects: the dataset is roughly four times
harder (80 classes versus 20, with a long tail of small and clustered
objects), the training schedule is half the duration, the backbone is
substantially smaller, and the augmentation pipeline is minimal.

The central finding of the project is methodological rather than numerical.
The first four training runs (V1 through V4), spanning two different
backbones and a deliberately wide sweep of learning rates and gradient-clip
thresholds, all produced exactly $\mathrm{mAP}@0.5 = 0$ on the test set
despite training and validation losses that decreased monotonically and
plateaued at values an observer would naturally interpret as evidence of a
learning model. The cause was not a poor hyper-parameter choice or a
fundamental limitation of the YOLO formulation, but two independent
implementation bugs:

\begin{itemize}
    \item A loss function that indexed the per-cell prediction tensor with
        positions appropriate for the Pascal VOC layout ($C=20$) while
        receiving COCO targets ($C=80$). Of $1\,062$ ground-truth object
        cells in a sample of 200 training images, the buggy
        ``exists\_box'' mask was set for only three of them (the elephants).
        The model was gradient-blind to the bounding-box slots it was
        supposed to predict, and the entire localisation signal evaporated.
    \item A post-processing function that divided predicted widths and
        heights by $S=7$ before returning them, even though the dataset and
        loss had agreed to express $w, h$ directly in image-relative units.
        Even a perfectly trained model would have produced boxes seven times
        smaller than ground truth, with a maximum achievable IoU of
        $\sim\!0.02$---well below the $0.5$ threshold of $\mathrm{mAP}@0.5$.
\end{itemize}

\noindent Each bug, in isolation, is sufficient to produce
$\mathrm{mAP}@0.5 = 0$ at evaluation time. Both were invisible to the
pre-existing unit-test suite, which checked tensor shapes and value ranges
but not quantitative coordinate or value invariants. Both were detected
within a single audit pass once the question shifted from
``why are the hyper-parameters wrong?'' to ``is the loss the loss?''.

The pedagogical takeaway, broader than the specific bugs, is that a
detection pipeline contains many silent failure modes whose individual
symptoms are too subtle to diagnose by training-curve inspection alone.
A converging loss curve is consistent both with a model learning what we
want it to learn and with a model learning an arithmetically valid quantity
that is not what we want. The distinction can be settled cheaply with
round-trip tests on each coordinate transformation and with loss tests that
exercise the production class count rather than only the count for which
the test was originally written. We have added both kinds of test alongside
the fixes; the project would have arrived at $\mathrm{mAP}@0.5 > 0$ within
the first training run if they had existed from the start.

Whether the resulting detector would have been competitive with modern
single-stage detectors is a question the architecture of YOLOv1 was never
designed to answer. The improvements that close most of the residual gap to
its descendants---anchor boxes, multi-scale feature pyramids, mosaic
augmentation, fully-convolutional heads---are by now the standard fare of
the YOLO family from \citet{redmon2018yolov3} onwards, and they constitute
the natural continuation of this work.
